This chapter documents the enhancements that would be required to transition
the platform from its current Docker Compose deployment to a
production-grade banking environment.

% ═══════════════════════════════════════════════════════════
\section{Container Orchestration with Kubernetes}
\label{sec:kubernetes}
% ═══════════════════════════════════════════════════════════

The current deployment runs 20 Docker Compose services on a single host.
This is sufficient for development and demonstration, but a production
deployment requires horizontal scaling, automatic failover, and
resource-aware scheduling across multiple nodes.

Kubernetes would replace Docker Compose as the orchestration layer. Each
streaming service (enrichment, scoring, decision) would run as a Deployment
with configurable replica counts, allowing the system to scale horizontally
under increased transaction volume. Health checks already exposed by the
Prometheus metrics endpoints can serve as Kubernetes liveness and readiness
probes with no code changes. The model volume mount --- currently a shared
Docker volume --- would migrate to a Persistent Volume Claim, preserving
the retrain-without-rebuild pattern established in
Chapter~\ref{ch:mlops}.

RedPanda, Redis, and PostgreSQL would each run as StatefulSets with
persistent storage, or be replaced by managed cloud equivalents depending
on the bank's infrastructure policy.


% ═══════════════════════════════════════════════════════════
\section{Continuous Integration and Deployment}
\label{sec:cicd}
% ═══════════════════════════════════════════════════════════

The project currently has no automated CI/CD pipeline. Code changes are
tested manually and deployed by rebuilding Docker images. A production
pipeline would introduce three stages:

\begin{enumerate}
    \item \textbf{Continuous Integration} --- on every commit, run unit
          tests for the three core modules (EnrichmentCore, ScoringCore,
          DecisionCore), validate that the feature computation produces
          the expected 30-dimensional output, and lint the codebase.
    \item \textbf{Image Build} --- on merge to the main branch, build and
          tag Docker images for each service with the commit hash,
          ensuring traceability between deployed code and source control.
    \item \textbf{Staged Deployment} --- deploy to a staging environment
          first, run the E2E test script against the staging pipeline,
          and promote to production only after validation passes.
\end{enumerate}

The core extraction pattern (Chapter~\ref{ch:streaming}) simplifies CI
significantly: the pure core modules can be unit-tested without Redis,
PostgreSQL, or RedPanda dependencies, covering the majority of the
business logic with fast, isolated tests.


% ═══════════════════════════════════════════════════════════
\section{Secrets Management}
\label{sec:secrets}
% ═══════════════════════════════════════════════════════════

Database credentials, API keys, and the Airflow REST API password are
currently stored in a \texttt{.env} file excluded from version control
via \texttt{.gitignore}. This is adequate for a single-developer
environment but insufficient for production, where secrets must be
rotated, audited, and never stored on disk in plaintext.

A secrets management solution such as HashiCorp Vault or cloud-native
equivalents (AWS Secrets Manager, Azure Key Vault) would provide
encrypted storage, automatic rotation, and access logging. Docker Compose
secrets or Kubernetes Secrets would inject credentials at runtime,
eliminating environment files entirely.


% ═══════════════════════════════════════════════════════════
\section{Load Testing}
\label{sec:load-testing}
% ═══════════════════════════════════════════════════════════

The system has been validated with the synthetic dataset of 243,000 events
replayed at accelerated speed, but it has not been tested under sustained
production-equivalent load. A load testing phase would establish the
throughput ceiling of each service, identify the first bottleneck in the
pipeline, and determine the point at which consumer lag begins to
accumulate.

The existing event simulator can serve as the load generator with minimal
modification: adjusting the inter-event delay to simulate peak branch
hours across 150 branches. The Prometheus instrumentation
(Section~\ref{sec:observability}) already captures the metrics needed to
evaluate load test results --- throughput, p95 latency, and consumer lag
--- without additional tooling.


% ═══════════════════════════════════════════════════════════
\section{Data Retention and Archival}
\label{sec:data-retention}
% ═══════════════════════════════════════════════════════════

The current PostgreSQL schema retains all transactions and alerts
indefinitely. In a production deployment, regulatory requirements and
storage constraints necessitate a data retention policy. Tunisian banking
regulations require that transaction records be preserved for a minimum
period, after which they may be archived to cold storage or purged.

A retention strategy would partition the \texttt{transactions} and
\texttt{alerts} tables by month using PostgreSQL's native table
partitioning, enabling efficient archival of older partitions to
compressed cold storage while keeping recent data queryable at full
speed. The batch profile rebuild
(Chapter~\ref{ch:streaming}) would be scoped to the retention window
rather than the full transaction history, reducing nightly computation
time as the dataset grows.